% 三、模型假设；四、符号说明
\section{模型假设}
\begin{enumerate}
\item \textbf{时间离散与段内恒定}。将附件中的功率序列按 10\,min 粒度离散，一天共 $T=144$ 个时段；每一时段内负载、光伏和储能充放电功率视为恒定，功率 $P$（kW）对应的时段电量为 $P/6$（kWh）。
\item \textbf{当前状态可观测}。执行进入时段 $t$ 后，当段实际负载、实际光伏出力、已经揭示的交易价格以及段初储电量均可获得，储能设备能够按照同一时间粒度完成当段控制；尚未实现的未来负载、光伏和价格不作为实际执行信息。
\item \textbf{信息使用满足因果性}。任一决策时刻的预测、残差场景和参数估计仅使用该时刻以前已经完成的历史观测以及已经正式发布的预报；问题 3 中的 0:00、6:00、12:00、18:00 光伏预报仅在对应发布时间之后进入模型。
\item \textbf{储能效率在研究期内恒定}。按题设将充电效率和放电效率均取 $\eta=0.9$，充电交流侧电量 $C_t$ 对库存增加 $\eta C_t$，放电交流侧电量 $D_t$ 对库存消耗 $D_t/\eta$；题目未提供参数的自放电、电池寿命衰减和启停成本不纳入模型。
\item \textbf{不考虑反向售电收益}。题目仅给出向外网购电和紧急购电的价格机制，未规定微网向外网售电的渠道及价格，因此超过负载和可充电需求的光伏或已购电量允许弃置，且不产生售电收入。
\item \textbf{问题 2--4 按跨日连续库存处理}。除问题 1 按照题目要求施加 $S_0=S_{144}$ 的周期约束外，问题 2--4 不要求每日回到固定库存，当日 24:00 的储电量作为下一日 0:00 的初始状态。为与当前计算结果保持一致，正式评分区间取 2 月 1 日至 12 月 31 日，并在评分期初统一以 6000\,kWh 初始化；1 月数据仅用于预测、场景和参数的因果预热。
\end{enumerate}

紧急购电仅用于补足当段供能缺口，这一处理直接对应题目对紧急购电触发条件的规定，不再作为独立经验假设。同一 10\,min 时段内的充、放电互斥则可由后文互斥消去引理得到，因此无需额外引入二元变量。

\section{符号说明}
记号约定：帽记号 $\hat x$ 表示预测量，上标 $\mathrm{act}$ 表示实际观测，$[z]^+=\max(z,0)$；凸函数 $f$ 的次微分记 $\partial f$。主要符号见表~\ref{tab:symbols}，均在首次出现处另有说明。

\begin{table}[!htbp]\centering\small
\caption{主要符号及适用范围}\label{tab:symbols}
\renewcommand{\arraystretch}{1.08}
\begin{tabularx}{\textwidth}{@{}lXcc@{}}\toprule
\tabheadrow 符号 & 含义 & 单位 & 适用 \\ \midrule
$t,\,T$ & 10 分钟段编号与一天段数，$T=144$ & --- & 一至四 \\
$L_t,\,V_t$ & 段 $t$ 负载、光伏电量 & kWh & 一至四 \\
$p_t$ & 段 $t$ 交易电价（问 1--3 为附件 1 常数曲线） & 元/kWh & 一至四 \\
$G_t,\,G^0_t$ & 购电量；0:00 冻结的计划购电量 & kWh & 一至四 \\
$C_t,\,D_t$ & 储能充、放电量（交流侧） & kWh & 一至四 \\
$S_t$ & 段末储电量，$S_t\in[1200,10800]$ & kWh & 一至四 \\
$\eta,\,\bar P$ & 充放电效率 $0.9$；每段充放电量上限 $5000/6$ & ---, kWh & 一至四 \\
$\psi(x)$ & 库存增量 $x$ 的交流侧电量：$x\ge0$ 为 $x/\eta$，$x<0$ 为 $\eta x$ & kWh & 一至四 \\
$F_t(e)$ & 问 1 确定性剩余最低费用（价值函数） & 元 & 一 \\
$E_t,\,W_t$ & 紧急购电量（5 倍电价）；富余电量 & kWh & 二至四 \\
$M$ & 配对残差场景数，$M=30$ & --- & 二至四 \\
$\omega$ & 场景编号，$\omega=1,\dots,M$ & --- & 二至四 \\
$H_{t,\omega},\,\bar H_t$ & 场景 $\omega$ 未来费用函数；其等权平均 & 元 & 二至四 \\
$R_t$ & 动态保留水平：缺口时不再放电的库存下限 & kWh & 二至四 \\
$V(e)$ & 次日价值函数（次日自由购电的最低费用） & 元 & 二至四 \\
$v$ & 线性续存费率（单日时域配置的末端项系数） & 元/kWh & 二至四 \\
$K_L,\,K_P$ & 负载、光伏历史窗参数（消融配置用） & --- & 二至四 \\
$\hat B_d,\,\hat s_{d,t},\,\beta_d$ & 预测日电量、归一化日内形状、类型切换率 & kWh, ---, --- & 二至四 \\
$\tau,\,q_\tau$ & 发布时刻（0/6/12/18 时）；当天剩余段数 $144-6\tau$ & ---, --- & 三、四 \\
$G^a_t$ & 调整购电量：重优化后对段 $t$ 提交的购电量 & kWh & 三、四 \\
$\Delta^+_t,\,\Delta^-_t$ & 调增、调减量，$G^a=G^0+\Delta^+-\Delta^-$ & kWh & 三、四 \\
$w$ & 0:00 正式预报与历史预测的组合权重 & --- & 三、四 \\
$\widetilde G_{\omega,u}$ & 跨午夜次日暂拟购电（按场景追索，不提交） & kWh & 三、四 \\
$\Gamma_\tau$ & 重优化目标中当天已定购电费与调整费部分 & 元 & 三、四 \\
$\ell_d,\,\chi_{d,t}$ & 日价格水平；归一化日内价格形状 & 元/kWh, --- & 四 \\
$K_p$ & 电价历史窗长度（35 个完整日） & --- & 四 \\
$\lambda$ & 价格同比缩放因子（命题~\ref{prop:scale}） & --- & 四 \\
$K_{\tau,\omega}(e)$ & 场景 $\omega$ 从午夜库存 $e$ 起的跨日续存费用 & 元 & 二至四 \\
\bottomrule\end{tabularx}\end{table}

标签 0:10 表示段 $[0{:}00,0{:}10)$，0:00+1 表示当天最后一段（24:00）。附件 3 的整点预报先插值到段末再转电量。四小时汇总表中充、放电量可同时非零，指该四小时内不同段分别发生充电与放电，同一 10 分钟段内充放电互斥。
